\chapter{张量、外代数与向量丛初步}
\label{ch:06-tensors}
\ChapterPlan{周2.1至2.3（本地页码待核对）；梁上册第二版2.4及第2章习题14、15、22；外代数与丛构造另作补充}{正文及证明初稿；附8题解答与原创图，尚未独立数学审阅}

第1章把多重线性形式作为预备，第4章构造了切丛与余切丛。现在从泛性质统一这些对象，区分张量、张量分量和张量场，并说明不同局部表达怎样拼成全局对象。本章只使用有限维实线性代数；无限维完备化、一般模的张量积和主丛理论不在这里展开。

\section{张量积的泛性质}
双线性映射$b:V\times W\to Z$通常不是乘积空间上的线性映射：同时缩放两个输入会产生二次因子。张量积将“分别线性”改写成某个新空间上的一次线性，从而让定义不依赖选基。
\begin{definition}[张量积的泛性质]
向量空间$V,W$的张量积（tensor product）是一对$(V\otimes W,\tau)$，其中$\tau:V\times W\to V\otimes W$双线性，并满足：对任意向量空间$Z$及双线性$b:V\times W\to Z$，有唯一线性映射$\widetilde b:V\otimes W\to Z$使$b=\widetilde b\circ\tau$。写$\tau(v,w)=v\otimes w$，称为纯张量（pure tensor）。
\end{definition}
\begin{theorem}[张量积的构造、唯一性与基]\label{thm:tensor-universal}
张量积存在，且在保持$\tau$的线性同构意义下唯一。若$e_1,\ldots,e_m$与$f_1,\ldots,f_n$分别为$V,W$的基，则$e_i\otimes f_j$为$V\otimes W$的基。因此$\dim(V\otimes W)=mn$，每个张量是有限个纯张量之和。
\end{theorem}
\begin{proof}
令$F$为以集合$V\times W$为形式基的自由向量空间，即元素为有限和$\sum a_{v,w}[v,w]$。取子空间$R$，由如下元素张成：
\begin{align*}
 &[v+v',w]-[v,w]-[v',w],\quad [cv,w]-c[v,w],\\
 &[v,w+w']-[v,w]-[v,w'],\quad [v,cw]-c[v,w].
\end{align*}
定义$V\otimes W=F/R$，$v\otimes w=[v,w]+R$。这些关系保证$\tau$双线性。给定$b$，按形式基定义$L([v,w])=b(v,w)$并线性延拓。双线性恰好保证$L$在$R$上为零，因此下降为$\widetilde b$。由于纯张量张成商空间，$\widetilde b$在它们上的值已经确定，故唯一。

若$(Q,\sigma)$也有此性质，分别得到$A:V\otimes W\to Q$、$B:Q\to V\otimes W$，将纯张量送到相应的$\sigma(v,w)$并送回。复合$BA$与恒等映射在所有纯张量上相同，故由唯一性相等；$AB$同理。这证明同构及其唯一性。

用基展开$v,w$，可见$e_i\otimes f_j$张成。为证独立性，固定$(a,b)$，双线性函数$(v,w)\mapsto e^a(v)f^b(w)$下降为商空间上的线性泛函。作用于$\sum c^{ij}e_i\otimes f_j=0$便得到$c^{ab}=0$。对所有$(a,b)$如此，故它们为基。零维因子时相应基为空，张量积为零空间，维数公式仍成立。
\end{proof}
\begin{figure}[htbp]
\centering\DGDiagram{tensor-universal}

\caption{原创示意：双线性映射$b$唯一经由$\tau$分解成线性映射$\widetilde b$。这里$V\times W$表示双线性映射的输入集合，箭头$\tau,b$按分别线性理解，不要求它们对乘积向量空间同时线性；虚线箭头表示由泛性质唯一确定的线性映射。}
\label{fig:tensor-universal}
\end{figure}
\begin{proposition}[映射、结合与交换同构]\label{prop:tensor-maps}
线性映射$A:V\to V'$、$B:W\to W'$唯一诱导$A\otimes B$，满足
\[
 (A\otimes B)(v\otimes w)=Av\otimes Bw.
\]
张量积具有自然同构
\[
 (U\otimes V)\otimes W\simeq U\otimes(V\otimes W),\quad
 V\otimes W\simeq W\otimes V,\quad \R\otimes V\simeq V.
\]
它们分别由去除括号、交换因子、标量乘法确定。诱导映射保持恒等与复合。
\end{proposition}
\begin{proof}
$(v,w)\mapsto Av\otimes Bw$双线性，应用泛性质即得诱导映射。两个复合在纯张量上的值都是相应线性映射复合后的张量积，故相等。
对结合同构，映射$(u,v,w)\mapsto u\otimes(v\otimes w)$逐槽线性；先按前两个变量使用泛性质，再对其结果与$w$使用泛性质，得到所需线性映射。反向按相反括号构造，两个复合在三因子纯张量上为恒等，因它们张成而为互逆。交换映射由$(v,w)\mapsto w\otimes v$给出，交换两次为恒等。最后$c\otimes v\mapsto cv$与$v\mapsto1\otimes v$互逆。
\end{proof}
由此可不写括号地使用多因子张量积；任意逐槽线性映射都唯一经由它分解。可用相同的形式符号商空间构造验证多因子泛性质，或按上述双因子过程归纳。
\begin{example}[和张量不一定为纯张量]\label{ex:tensor-rank}
在$\R^2\otimes\R^2$中，$T=e_1\otimes e_1+e_2\otimes e_2$不是纯张量。若$T=v\otimes w$，其分量矩阵为$vw^{\mathsf T}$，秩至多为一；$T$的矩阵是单位矩阵，秩为二，矛盾。相反，若非零分量矩阵秩一，取一个非零列$v$，所有列均为其倍数，便能将矩阵写成$vw^{\mathsf T}$，因此相应张量为纯张量。纯张量构成的集合通常不是线性子空间。
\end{example}
\begin{proposition}[对偶与线性算子]\label{prop:tensor-hom}
有限维空间有不依赖选基的同构
\[
 V\simeq V^{**},\qquad
 V^*\otimes W\simeq\operatorname{Hom}(V,W),\qquad
 (V\otimes W)^*\simeq V^*\otimes W^*.
\]
第一映射是$v\mapsto(\lambda\mapsto\lambda(v))$；第二映射把$\lambda\otimes w$送到$v\mapsto\lambda(v)w$；第三同构由对偶张量在纯张量上的乘积配对确定。
\end{proposition}
\begin{proof}
若$v\ne0$，某个基分量非零，相应对偶函数在$v$上不为零，故求值映射单射。两边同维，所以是同构；定义没有选择内积。
第二映射由双线性性良定义。给定$A$，其逆像可写为$\sum_i e^i\otimes Ae_i$，因为该张量作用于$v=v^ie_i$得到$Av$。而$e^i\otimes f_j$对应通常的矩阵单位，线性独立，所以是同构；逆像表达虽用了基，同构本身并未使用基。
最后定义$(\lambda\otimes\mu)(v\otimes w)=\lambda(v)\mu(w)$。基$e^i\otimes f^j$正好映为$e_i\otimes f_j$的对偶基，故得到第三个同构。这些结论使用有限维假设，不据此断言任意无限维线性算子都是有限秩张量。
\end{proof}
交换同构的存在不等于两个张量本来相等。在$V\otimes V$中，一般$v\otimes w\ne w\otimes v$。本章的普通张量交换没有符号；交替代数中的符号来自另加的商关系，见后节。习题6.1检验泛性质、维数与可分解性。


\Needspace{22\baselineskip}
\section{张量型别与缩并}
\begin{definition}[张量的型别与分量]
对有限维$V$，定义
\[
 T^r_s(V)=V^{\otimes r}\otimes(V^*)^{\otimes s},\qquad T^0_0(V)=\R.
\]
其元素称$(r,s)$型张量（tensor of type $(r,s)$）。由上一节的求值同构，可将它等价视为$(V^*)^r\times V^s$上的逐槽线性实函数：上标槽输入协向量，下标槽输入向量。基下有唯一展开
\[
 T=T^{i_1\cdots i_r}{}_{j_1\cdots j_s}
 e_{i_1}\otimes\cdots\otimes e_{i_r}\otimes
 e^{j_1}\otimes\cdots\otimes e^{j_s}.
\]
分量通过在相应槽中代入对偶基和原基得到；维数为$(\dim V)^{r+s}$（零次张量约定维数一）。
\end{definition}
等价性的证明可逐个基张量进行：上述基张量在基输入元组上给出Kronecker符号乘积，这些乘积构成全部逐槽线性形式的基。因此不是把“有若干指标的数组”直接当作无坐标定义。
\begin{theorem}[张量换基律]\label{thm:tensor-transform}
若$e'_a=e_iP^i{}_a$，$P$可逆，则
\[
 T'^{a_1\cdots a_r}{}_{b_1\cdots b_s}
 =(P^{-1})^{a_1}{}_{i_1}\cdots(P^{-1})^{a_r}{}_{i_r}
 P^{j_1}{}_{b_1}\cdots P^{j_s}{}_{b_s}
 T^{i_1\cdots i_r}{}_{j_1\cdots j_s}.
\]
逆变分量乘$P^{-1}$，协变分量乘$P$。对于流形坐标变换$y=y(x)$，记$J^a{}_i=\partial y^a/\partial x^i$，则切空间新基对应$P=J^{-1}$。
\end{theorem}
\begin{proof}
由第1章对偶换基式，$e'^a=(P^{-1})^a{}_i e^i$。分量$T'$是在张量的上槽依次输入$e'^{a_1},\ldots,e'^{a_r}$、下槽输入$e'_{b_1},\ldots,e'_{b_s}$后求值。逐槽线性使各$P$或$P^{-1}$因子提出，给出公式。坐标情形由$\partial/\partial y^a=(J^{-1})^i{}_a\partial/\partial x^i$得到。
\end{proof}
\begin{definition}[张量积与缩并]
两个混合型张量的积先作普通张量积，再将所有$V$因子保持原顺序排到前面、所有$V^*$因子保持原顺序排到后面，得到型别相加的张量；这里的重新排序通过上一节的交换同构实现，无额外符号。
对$r,s\geq1$，缩并（contraction）$C^a_b:T^r_s(V)\to T^{r-1}_{s-1}(V)$在纯张量上定义为：将第$b$个协向量作用于第$a$个向量，乘所得标量，再删去这两个因子，其余因子顺序保留。
\end{definition}
\begin{theorem}[缩并的良定义与基无关性]\label{thm:contraction}
缩并是与基无关的线性映射。在分量上，它把一个上标和一个下标置成同一指标并求和。例如
\[
 (C^1_1T)^{i_2\cdots i_r}{}_{j_2\cdots j_s}
 =\sum_{a=1}^{\dim V}T^{a i_2\cdots i_r}{}_{a j_2\cdots j_s}.
\]
特别地，$V\otimes V^*\simeq\operatorname{End}(V)$下缩并即迹，且$C(v\otimes\lambda)=\lambda(v)$。不指定额外双线性形式时，没有这种自然的两个上标或两个下标缩并。
\end{theorem}
\begin{proof}
对各向量与协向量，定义中的输出是逐槽线性的：被删去的两个槽通过双线性配对$\lambda(v)$进入标量，其余槽线性进入剩余张量。因此多因子泛性质给出唯一线性映射，定义本身没有用基。
将$T$按基展开，配对$e^b(e_a)=\delta^b_a$只留下两个指标相等的项，即得分量公式。换基时，被缩并两槽出现的逆矩阵与矩阵相消，剩余槽恰按低一型别的变换律变换。对算子$A$，其张量为$A^i{}_j e_i\otimes e^j$，缩并为$\sum_i A^i{}_i$，即迹。最后自然配对只接受一个$V$和一个$V^*$，两个同型槽之间尚无指定的配对，不能自行求和当成不变量。习题6.2专门核对$(2,1)$型的换基求和。
\end{proof}
\begin{example}[双线性形式、升降指标与迹]
双线性形式$g\in T^0_2(V)$在基下为$g_{ij}$。若$g$非退化，$g^\flat:V\to V^*$、$g^\flat(v)=g(v,\cdot)$为同构；对称时其逆用矩阵$(g^{ij})=(g_{ij})^{-1}$表示。于是$v_i=g_{ij}v^j$、$\lambda^i=g^{ij}\lambda_j$。这些升降指标操作依赖选定的$g$，并非张量定义自动附带的等同。习题6.3给出完整的非退化性论证。
在$V=\R^2$中，$B=e^1\otimes e^1+e^2\otimes e^2$为协变二阶张量。若$e'_1=2e_1,e'_2=e_2$，其矩阵从$I$变为$\operatorname{diag}(4,1)$，对角和从$2$变为$5$，说明直接相加$B_{11}+B_{22}$不是自然缩并。若另指定$g$，$g^{ij}B_{ij}$才是缩并得到的标量。
\end{example}
\begin{proposition}[协变拉回与型别边界]\label{prop:covariant-pullback}
线性$A:V\to W$对协变张量定义
\[
 (A^*T)(v_1,\ldots,v_s)=T(Av_1,\ldots,Av_s).
\]
该拉回保持张量积及复合。若$A$可逆，则也可将任意混合型张量拉回：向量因子用$A^{-1}$，协向量因子用$A^*$。一般非可逆$A$不由这些数据给出任意混合型张量的同型拉回。
\end{proposition}
\begin{proof}
协变定义逐槽线性，因此是$T^0_s(V)$的元素；代入纯张量可检验积与复合公式，任意张量再按基线性延拓。可逆时，在$r$个向量因子上取$A^{-1}$、$s$个对偶因子上取$A^*$的张量积，即得混合型定义。不可逆时$A^{-1}$不存在，因而不能把协向量的自然拉回无条件推广到所有型别。
\end{proof}
\begin{example}[普通偏导不自动成为张量]
在半轴$x>0$取$X=\partial_x$，换坐标$y=x^2$。新分量为$X^y=2\sqrt y$，所以$\partial_yX^y=1/\sqrt y$；旧坐标中$\partial_xX^x=0$。若“分量偏导”是$(1,1)$型张量，一维时上、下换元因子应相消，两者须相等，矛盾。联络将来要补偿这种换坐标时出现的额外项；这里不提前引入尚未定义的协变导数。
\end{example}


\Needspace{10\baselineskip}
\section{交替化、楔积与外代数}
张量积保留变量的次序。描述有向面积或体积时，还希望交换两个输入便变号；外代数将这一要求编码为商关系。为避免归一化混淆，先区分商空间中的外积与普通协变张量的交替投影。
\begin{definition}[张量代数与外代数]
令$V^{\otimes0}=\R$，张量代数为直和
\[
 \mathcal T(V)=\bigoplus_{k\geq0}V^{\otimes k},
\]
乘法为因子串接。直和中的每个元素只有有限个非零齐次分量。令$I$为由所有$v\otimes v$生成的双边理想，定义外代数（exterior algebra）$\Lambda V=\mathcal T(V)/I$，其$k$次齐次部分记为$\Lambda^kV$，乘法记为$\wedge$。
\end{definition}
\begin{theorem}[外代数的基与泛性质]\label{thm:exterior-universal}
若$\dim V=n$、基为$e_1,\ldots,e_n$，则$\Lambda^kV$的基为
\[
 e_{i_1}\wedge\cdots\wedge e_{i_k},\qquad i_1<\cdots<i_k.
\]
因此$\dim\Lambda^kV=\binom nk$，$k>n$时为零，$\Lambda^0V=\R$。任何交替的$k$重线性映射$b:V^k\to W$唯一分解为从$\Lambda^kV$到$W$的线性映射。外积结合，并对齐次元素满足
\[
 \xi\wedge\eta=(-1)^{pq}\eta\wedge\xi,
 \qquad\xi\in\Lambda^pV,\ \eta\in\Lambda^qV.
\]
\end{theorem}
\begin{proof}
生成理想的关系齐次，故商代数仍按次数分级。由$(v+w)\wedge(v+w)=0$及$v\wedge v=w\wedge w=0$，得到$v\wedge w=-w\wedge v$。于是任何基单项式可通过相邻交换排成递增指标；若指标重复，则移到相邻后为零。这证明所列单项式张成。

为证独立，对递增指标组$J=(j_1,\ldots,j_k)$定义行列式形式
\[
 b_J(v_1,\ldots,v_k)=\det(e^{j_a}(v_b))_{a,b=1}^k.
\]
其在$V^{\otimes k}$上的线性延拓对任何含相邻相同向量的纯张量为零，因此杀死理想$I$的$k$次部分，下降为$\Lambda^kV$上的线性泛函。在递增基单项式上取值为$\delta_{IJ}$，所以它们独立。这里若有非递增或重复的指标，由行列式的符号和重复列为零处理；并未预设待证的基结论。

同理，任意交替$b$的张量积线性延拓杀死相邻重复向量，故杀死$I$的相应齐次部分，给出因子分解；因纯外积张成而唯一。结合律从商代数的乘法继承。将$p$个一次因子与$q$个一次因子交换共需$pq$次相邻交换，因此在纯齐次单项式上得符号$(-1)^{pq}$，再由双线性性推广。
\end{proof}
对$V^*$应用上一定理，通过行列式求值将$\Lambda^k(V^*)$与第1章的交替$k$形式空间自然同构。下文仍写$\Lambda^kV^*$，不再区分这两个已经识别的模型。这个识别使用的是完整行列式，不能漏掉阶乘。
\begin{definition}[对称化与交替化]
对$k$重协变张量$T$，定义
\begin{align*}
 (\operatorname{Sym}_kT)(v_1,\ldots,v_k)
 &=\frac1{k!}\sum_{\sigma\in S_k}T(v_{\sigma(1)},\ldots,v_{\sigma(k)}),\\
 (\operatorname{Alt}_kT)(v_1,\ldots,v_k)
 &=\frac1{k!}\sum_{\sigma\in S_k}\operatorname{sgn}(\sigma)
 T(v_{\sigma(1)},\ldots,v_{\sigma(k)}).
\end{align*}
零次情形二者均为恒等。此处二者是普通张量空间中的线性算子，交替化的像与上述交替形式空间相同。
\end{definition}
\begin{proposition}[投影与外积的归一化]\label{prop:alt-normalization}
$\operatorname{Sym}_k$与$\operatorname{Alt}_k$分别投影到对称和交替张量子空间，均满足幂等性。若$\alpha\in\Lambda^pV^*$、$\beta\in\Lambda^qV^*$，则在形式模型中
\[
 \alpha\wedge\beta=\frac{(p+q)!}{p!q!}
 \operatorname{Alt}_{p+q}(\alpha\otimes\beta).
\]
特别地，$\lambda\wedge\mu=\lambda\otimes\mu-\mu\otimes\lambda$，而
$\operatorname{Alt}_2(\lambda\otimes\mu)=\tfrac12\lambda\wedge\mu$。
\end{proposition}
\begin{proof}
将一个置换作用于求和式，相当于重新编号全部置换；不带符号时和保持，带符号时提出该置换的符号。因此像分别对称和交替。若输入已经对称，每个对称化求和项相同；若输入已经交替，每个带符号求和项也相同。于是相应算子在像上为恒等，得到幂等性。

对于外积，先对基外积验证。若$\alpha=e^{i_1}\wedge\cdots\wedge e^{i_p}$、$\beta=e^{j_1}\wedge\cdots\wedge e^{j_q}$，将两个行列式展开再交替化，结果为全部$p+q$个协向量的行列式乘$p!q!/(p+q)!$；因组内置换均由各自的交替性抵消符号，真正不同的项恰是保持两组内部次序的洗牌。乘以所列系数后得到合并的外积。按基双线性延拓即得一般结论。这也与第1章定理\ref{thm:alternating}的$1/(p!q!)$全排列求和完全一致。
\end{proof}
对二阶协变张量，有$T=\operatorname{Sym}_2T+\operatorname{Alt}_2T$，两子空间交为零；但$k\geq3$时不能据此声称仅对称与全交替两部分已穷尽全部张量。习题6.5检验归一化及偶次外积的边界。
\begin{proposition}[外积拉回与内插]\label{prop:exterior-pullback}
对线性$A:V\to W$，协变拉回满足$A^*(\alpha\wedge\beta)=A^*\alpha\wedge A^*\beta$。对$v\in V$定义内插（interior product）
\[
 (\iota_v\alpha)(v_2,\ldots,v_p)=\alpha(v,v_2,\ldots,v_p),
 \qquad\iota_v c=0\quad(c\in\R).
\]
则$\iota_v^2=0$且
\[
 \iota_v(\alpha\wedge\beta)
 =(\iota_v\alpha)\wedge\beta+(-1)^p\alpha\wedge\iota_v\beta.
\]
\end{proposition}
\begin{proof}
拉回与交替化逐项交换，且保持普通张量积；应用上一命题的公式便得外积拉回式。内插后仍交替；连续插入两次相同$v$，交替性给出零。
若$\alpha=\lambda_1\wedge\cdots\wedge\lambda_p$，按行列式的第一列展开可得
\[
 \iota_v\alpha=\sum_{a=1}^p(-1)^{a-1}\lambda_a(v)
 \lambda_1\wedge\cdots\wedge\widehat\lambda_a\wedge\cdots\wedge\lambda_p.
\]
对$\alpha\wedge\beta$中的全部一次因子套用此式，插入落在前$p$因子的项组成第一项，落在后$q$因子的项先跨过$p$个因子而多出$(-1)^p$，组成第二项。纯外积张成，故线性延拓成立；零次情形按定义也相容。帽号表示删去相应因子。
\end{proof}
\begin{example}[偶次形式的平方未必为零]
奇数次齐次形式$\alpha$满足$\alpha\wedge\alpha=-\alpha\wedge\alpha$，在实数域上即为零。偶数次则未必：在$\R^4$上取
\[
 \omega=e^1\wedge e^2+e^3\wedge e^4,
 \qquad \omega\wedge\omega=2e^1\wedge e^2\wedge e^3\wedge e^4\ne0.
\]
因此$\omega$不能写成两个一次形式的单个外积，后者平方因出现重复因子而必为零。不要把“所有张量是纯张量”或“所有二形式都可分解”混作定义的一部分。
\end{example}
顶次外幂$\det V=\Lambda^nV$是一维空间。若$A:V\to V$线性，将$Ae_1,\ldots,Ae_n$按基展开并利用重复因子为零，得到
\[
 (\Lambda^nA)(e_1\wedge\cdots\wedge e_n)
 =\det(A)\,e_1\wedge\cdots\wedge e_n.
\]
这给出行列式的无坐标解释。它将来连接体积形式、定向与积分；本章不由这一点预先宣称所有流形都存在全局非零体积形式。


\Needspace{10\baselineskip}
\section{向量丛、截面及过渡函数}
逐点张量还需具备随基点光滑变化的规则。向量丛（vector bundle）用局部乘积和纤维中的线性过渡同时表达这两层要求；第4章的切丛、余切丛正是特例。
\begin{definition}[光滑实向量丛]
秩$k$的光滑实向量丛为光滑投影$\pi:E\to M$，每条纤维$E_p=\pi^{-1}(p)$具有$k$维实向量空间结构，且每点附近存在微分同胚
\[
 \phi_\alpha:\pi^{-1}(U_\alpha)\longrightarrow U_\alpha\times\R^k
\]
保持基点并在每个纤维上线性。称$\phi_\alpha$为局部平凡化，总空间$E$按全书约定是Hausdorff、第二可数流形。秩零时允许零空间纤维。光滑截面为满足$\pi\circ s=\operatorname{id}_M$的光滑映射$s:M\to E$。
\end{definition}
\Needspace{5\baselineskip}
局部标架$e_{\alpha1},\ldots,e_{\alpha k}$由$e_{\alpha i}(p)=\phi_\alpha^{-1}(p,\varepsilon_i)$给出，其中$\varepsilon_i$是$\R^k$标准基。截面$s=e_{\alpha i}s_\alpha^i$光滑当且仅当全部系数光滑，因为其局部表达为$p\mapsto(p,s_\alpha(p))$。
\begin{theorem}[过渡函数与拼接构造]\label{thm:bundle-gluing}
局部平凡化满足
\[
 \phi_\beta\phi_\alpha^{-1}(p,v)=(p,g_{\beta\alpha}(p)v),
 \qquad g_{\beta\alpha}:U_\alpha\cap U_\beta\to\operatorname{GL}(k,\R)
\]
光滑，且在重叠处满足
\[
 g_{\alpha\alpha}=I,\quad g_{\alpha\beta}=g_{\beta\alpha}^{-1},\quad
 g_{\gamma\alpha}=g_{\gamma\beta}g_{\beta\alpha}.
\]
反之，在光滑流形$M$的开覆盖上给定这些光滑矩阵函数及相容条件，就可构造一个向量丛，局部过渡恰为这些函数。固定这些局部识别时，所得丛在丛同构意义下唯一。
\end{theorem}
\begin{proof}
正向由纤维线性性得到矩阵形式；矩阵每列是光滑过渡映射作用于固定标准基的结果，所以系数光滑。恒等、逆及复合条件由$\phi_\gamma\phi_\alpha^{-1}=(\phi_\gamma\phi_\beta^{-1})(\phi_\beta\phi_\alpha^{-1})$得到。注意矩阵次序是从$\alpha$坐标出发，先到$\beta$，再到$\gamma$。

反向取不交并$\bigsqcup_\alpha U_\alpha\times\R^k$，规定
\[
 (p,v)_\alpha\sim(p,w)_\beta\quad\Longleftrightarrow\quad
 w=g_{\beta\alpha}(p)v.
\]
三个相容条件分别保证自反、对称、传递，因此得到商集合$E$及投影$\pi([p,v])=p$。每个固定$\alpha$的映射$j_\alpha:(p,v)\mapsto[p,v]_\alpha$单射。给$E$商拓扑。局部开集$O\subset U_\alpha\times\R^k$在任意另一块中的等价原像由光滑可逆过渡送来，仍开，所以$j_\alpha(O)$在商中开。因此$j_\alpha$为到$\pi^{-1}(U_\alpha)$的同胚，逆给出所需平凡化。

不同基点的两个总空间点由基空间的Hausdorff性分开；相同基点的不同点在同一平凡化中由不交欧氏邻域分开，因块像开，这些也是总空间开邻域。故$E$为Hausdorff。$M$的第二可数性给出该覆盖的可数子覆盖：对可数基中每个包含于某个$U_\alpha$的成员，选一个这样的$U_\alpha$即可。对应的可数个局部乘积各有可数基，它们的并是$E$的可数基。结合基空间坐标与纤维坐标，过渡光滑，故$E$为$\dim M+k$维光滑流形。

同一纤维中不同代表之间通过可逆线性$g$转换，所以向量空间结构一致，平凡化纤维线性。若另一丛有相同的局部识别，就将$[p,v]_\alpha$送到其$\alpha$平凡化的逆像；相容条件保证良定义，各块上为微分同胚，故得到全局丛同构。
\end{proof}
\begin{figure}[htbp]
\centering\DGDiagram{bundle-cocycle}

\caption{原创示意：固定三重交叠中的同一基点$p$，纤维坐标沿$\alpha\to\beta\to\gamma$的变换必须与直接$\alpha\to\gamma$相同，即$g_{\gamma\alpha}=g_{\gamma\beta}g_{\beta\alpha}$。图只画同一纤维的三种坐标表达，不是三条不同纤维。}
\label{fig:bundle-cocycle}
\end{figure}
\begin{proposition}[截面拼接、换平凡化与全局标架]\label{prop:bundle-sections}
局部光滑系数$f_\alpha:U_\alpha\to\R^k$给出全局光滑截面当且仅当$f_\beta=g_{\beta\alpha}f_\alpha$。若局部坐标另换为$v'_\alpha=h_\alpha v_\alpha$，其中$h_\alpha:U_\alpha\to\operatorname{GL}(k,\R)$光滑，则过渡矩阵变为
\[
 g'_{\beta\alpha}=h_\beta g_{\beta\alpha}h_\alpha^{-1}.
\]
$E$平凡（即与$M\times\R^k$丛同构）当且仅当存在$k$个全局光滑截面在每点组成纤维的基。
\end{proposition}
\begin{proof}
同一点的两种代表相等恰由拼接关系刻画，故得到截面的充要条件，光滑性可在每个局部乘积中检验。将$v_\alpha=h_\alpha^{-1}v'_\alpha$代入$v'_\beta=h_\beta g_{\beta\alpha}v_\alpha$得到换平凡化公式。
若有全局标架$s_1,\ldots,s_k$，定义$(p,a)\mapsto\sum_i a^is_i(p)$。局部标架下该映射的纤维矩阵光滑可逆，逆矩阵也光滑，故是全局丛同构。反之，把平凡丛的常值基截面经同构移回即得全局标架。秩零时为空标架，结论相同。
\end{proof}
截面空间$\Gamma(E)$是$C^\infty(M)$上的模，可逐点相加和乘函数。上一定理没有断言任意丛都有全局标架；局部线性代数不能替代全局拼接条件。
\Needspace{15\baselineskip}
\begin{example}[Möbius线丛]\label{ex:mobius-line}
在$\R\times\R$上作等同$(t,a)\sim(t+2\pi,-a)$，投影到$S^1=\R/2\pi\mathbb Z$，得到实线丛$L$。在长度小于$2\pi$的基空间角区间上选一个角提升，便得到局部乘积；角提升相差$2\pi m$时纤维坐标乘$(-1)^m$，故过渡光滑且满足相容条件。定理\ref{thm:bundle-gluing}保证这是光滑秩一丛。
若它有全局无零截面，拉回到实线便给出光滑函数$a(t)$满足$a(t+2\pi)=-a(t)$且处处非零。但连续函数从$a(0)$走到$-a(0)$必过零，矛盾。因此没有全局标架，$L$不平凡。这个证明只用中间值定理，无需提前使用特征类。
\end{example}
\begin{figure}[htbp]
\centering\DGDiagram{mobius-gluing}

\caption{原创拼接示意：Möbius线丛将$(0,a)$与$(2\pi,-a)$等同。蓝曲线画系数$a(t)=\cos(t/2)$，端点按反号关系拼接为全局截面，但它在$t=\pi$为零。两端橙色箭头表示纤维坐标方向的反向识别；矩形只是截取有限纤维范围的基本域，并非丛在三维空间中的等距嵌入。}
\label{fig:mobius-gluing}
\end{figure}
\begin{theorem}[丛的代数运算与张量场]\label{thm:tensor-bundles}
对同一底流形上的秩$m,n$丛$E,F$，逐点定义$E^*,E\oplus F,E\otimes F,\Lambda^kE$，可自然赋予向量丛结构，秩分别为$m,m+n,mn,\binom mk$。若原纤维坐标过渡为$g_{\beta\alpha},h_{\beta\alpha}$，新过渡分别为
\[
 g_{\beta\alpha}^{-\mathsf T},\quad
 \operatorname{diag}(g_{\beta\alpha},h_{\beta\alpha}),\quad
 g_{\beta\alpha}\otimes h_{\beta\alpha},\quad
 \Lambda^k g_{\beta\alpha}.
\]
特别地，
\[
 T^r_sM=(TM)^{\otimes r}\otimes(T^*M)^{\otimes s}
\]
是张量丛，其光滑截面称张量场。$\Lambda^kT^*M$的光滑截面称微分$k$形式。
\end{theorem}
\begin{proof}
在各原平凡化中选择对偶基、直和基、张量积基或递增外积基。坐标变换即所列矩阵：对偶坐标按逆转置变化；直和按块变化；张量积逐因子变化；外幂矩阵的系数是原矩阵相应小行列式。因逆矩阵光滑，其余系数为原系数的多项式，故全都光滑可逆。各运算保持复合，故满足相同的相容条件。例如$(AB)^{-\mathsf T}=A^{-\mathsf T}B^{-\mathsf T}$，顺序与纤维坐标复合一致；$\Lambda^k(AB)=(\Lambda^kA)(\Lambda^kB)$可在纯外积上验证。应用拼接定理得到丛结构及其唯一性，各点的秩由已证明的线性代数维数公式给出。

对$TM$使用坐标标架，便得到上一节的张量分量换元律；对$T^*M$得到一形式的规律。截面在每张卡中光滑等价于分量光滑，不同卡间按相同过渡相容，所以得到无坐标的张量场与微分形式定义。
\end{proof}
逐点张量积、缩并、外积及内插的分量为局部分量的有限和与乘积，因此将光滑截面送到光滑截面；基无关性保证局部结果拼接。对$F:M\to N$，协变场可用$\dd F_p$逐点拉回，分量为光滑函数与$DF$的有限乘积，故仍光滑。更一般地，拉回丛$F^*E=\{(p,e):F(p)=\pi(e)\}$在$F^{-1}(U_\alpha)$上有自然平凡化，过渡为$g_{\beta\alpha}\circ F$，拼接定理保证其丛结构。原截面$s$给出截面$p\mapsto(p,s(F(p)))$。这与把任意向量场直接拉回成源流形向量场是不同的目标空间。
\Needspace{10\baselineskip}
\begin{proposition}[常秩丛映射的核与像]\label{prop:bundle-constant-rank}
设$A:E\to F$为覆盖同一$M$的光滑丛映射（各纤维上线性），且各$A_p$秩恒为$r$。则$\ker A$与$\operatorname{im}A$分别为秩$m-r$与$r$的光滑子丛。若秩跳跃，不能直接得此结论。
\end{proposition}
\begin{proof}
在局部标架中$A$是光滑矩阵。选一个$r$阶可逆子块并重排行列，写作
\[
 A(p)=\begin{pmatrix}B&C\\D&E\end{pmatrix},\qquad B\text{可逆}.
\]
取可逆光滑矩阵
\[
 L=\begin{pmatrix}B^{-1}&0\\-DB^{-1}&I\end{pmatrix},\qquad
 R=\begin{pmatrix}I&-B^{-1}C\\0&I\end{pmatrix}.
\]
直接相乘得$LAR=\operatorname{diag}(I_r,E-DB^{-1}C)$。秩恒为$r$迫使右下块为零。因此在改造后的源、目标标架中，核由源的后$m-r$个基截面张成，像由目标的前$r$个基截面张成。这些局部线性子空间本身由点态核与像唯一确定，在重叠处相同；局部标架过渡光滑，故构成子丛。$r=0$时映射为零，结论直接成立；其余块为空的端点情况同样处理。
\end{proof}
第5章的光滑常秩分布因而可视为$TM$的子丛，其积分条件仍需李括号判别。下一章在$\Lambda^kT^*M$截面上建立外微分；本章不从逐点外代数运算推断$d$的存在，也不预设每个向量丛都有全局平凡化、度规或联络。


\Needspace{15\baselineskip}
\section{张量与向量丛习题及详细解答}
习题6.2、6.3、6.4分别改述梁书上册第二版第2章习题14、15、22（印刷页53，PDF物理页69），并增加必要的计算或边界说明。其余5题为原创补充。原题号、来源类型及正文核对范围见\texttt{references/EXERCISES.md}；解答均由本工程编写。

\Needspace{10\baselineskip}
\subsection*{习题6.1：泛性质与不可分解张量}
用泛性质证明$\operatorname{Hom}(V\otimes W,Z)$与双线性映射空间$\operatorname{Bil}(V,W;Z)$自然线性同构。取$V=W=\R^2$，证明$e_1\otimes e_1+e_2\otimes e_2$至少需要两个纯张量相加，并解释这为何不与“纯张量张成”矛盾。
\begin{proof}[解答]
将$L$送到$b_L(v,w)=L(v\otimes w)$，所得$b_L$双线性；反向由定理\ref{thm:tensor-universal}得到唯一$\widetilde b$。两者互逆，且加法、标量乘法都在生成元$v\otimes w$上逐项保持，故为线性同构。定义不选基，并且在输入、输出上线性复合与此对应相容，故称自然。
所给张量已写为两个纯张量的和，所以至多需要两个。若只需一个，则对应矩阵为$vw^{\mathsf T}$，秩至多为一，却应等于秩二的单位矩阵，矛盾。因此恰需两个。张成只断言每个元素可写成有限线性组合，既不保证只需一个生成元，也不保证这种分解唯一。
\end{proof}

\Needspace{12\baselineskip}
\subsection*{习题6.2：缩并与基的选择}
改述梁书第2章习题14：对$(2,1)$型张量$T$，证明缩并第一个上标与下标所得向量不依赖基。请直接写出两套基之间的求和计算，而不只引用正文的无坐标定义。
\begin{proof}[解答]
设$e'_a=e_iP^i{}_a$，缩并所得旧分量为$u^j=\sum_iT^{ij}{}_i$。新分量根据张量换基律满足
\begin{align*}
 u'^b&=\sum_aT'^{ab}{}_a\\
 &=\sum_{a,i,j,k}(P^{-1})^a{}_i(P^{-1})^b{}_jP^k{}_a T^{ij}{}_k\\
 &=\sum_{i,j,k}(P^{-1})^b{}_j\delta^k_iT^{ij}{}_k
 =\sum_j(P^{-1})^b{}_j u^j.
\end{align*}
故$u'=P^{-1}u$，两套分量代表同一向量，缩并与选基无关。
\end{proof}

\Needspace{10\baselineskip}
\subsection*{习题6.3：度规诱导的对偶同构}
改述梁书第2章习题15：设$g$为有限维实空间$V$上的对称非退化双线性形式，证明$g^\flat(v)=g(v,\cdot)$是$V\to V^*$的同构。给出逆映射的分量，说明正定性是否必要；若$g$退化又如何？
\begin{proof}[解答]
$g^\flat$的线性性来自$g$的第一槽线性。若$g^\flat(v)=0$，则$g(v,w)=0$对所有$w$成立，非退化性迫使$v=0$，故单射；有限维且$\dim V=\dim V^*$，所以满射。其矩阵在对称约定下为$(g_{ij})$，非退化等价于矩阵可逆；逆矩阵记为$(g^{ij})$，满足$g^{ij}g_{jk}=\delta^i_k$，因此$\lambda\mapsto g^{ij}\lambda_j e_i$为逆。
正定性没有用于证明，例如$\operatorname{diag}(1,-1)$同样给出同构；但若退化，则存在非零$v$使$g(v,w)=0$对所有$w$成立，故$g^\flat$有非零核，不再可逆。与$V\simeq V^{**}$不同，$V\simeq V^*$在这里依赖额外选择的$g$。
\end{proof}

\Needspace{11\baselineskip}
\subsection*{习题6.4：度规行列式不是标量}
改述梁书第2章习题22：设$g$为流形上的协变二阶张量场，在两套坐标$x,y$中的矩阵分别为$G_x,G_y$，令$K=\partial x/\partial y$。证明$G_y=K^{\mathsf T}G_xK$及$\det G_y=(\det K)^2\det G_x$。以平面欧氏度量及$y^1=2x^1,y^2=x^2$核查。
\begin{proof}[解答]
坐标基满足$\partial_{y^a}=K^i{}_a\partial_{x^i}$，所以
\[
 (G_y)_{ab}=g(\partial_{y^a},\partial_{y^b})
 =K^i{}_aK^j{}_b(G_x)_{ij}.
\]
这就是矩阵合同式；由行列式乘法性立即得到平方因子。特别地，分量行列式的数值一般依赖坐标，但非退化性和行列式的符号不变，因为$(\det K)^2>0$。
给定例中$G_x=I$，$K=\operatorname{diag}(1/2,1)$，故$G_y=\operatorname{diag}(1/4,1)$，行列式从$1$变为$1/4$。这是同一个几何张量的两种矩阵，不是几何度量改变；也不能因此把$\det G$本身当作全局标量函数。
\end{proof}

\Needspace{11\baselineskip}
\subsection*{习题6.5：交替化的系数与外积平方}
在$\R^4$上取$T=e^1\otimes e^2$与$\omega=e^1\wedge e^2+e^3\wedge e^4$。求$\operatorname{Alt}_2T$，并用其在$(e_1,e_2)$上的值检验系数。再求$\omega\wedge\omega$和$\iota_{e_1}\omega$，判断$\omega$能否是单个$\lambda\wedge\mu$。
\begin{proof}[解答]
交替投影为$\tfrac12(e^1\otimes e^2-e^2\otimes e^1)$，作用于$(e_1,e_2)$为$1/2$；而$e^1\wedge e^2$在该元组上为$1$，因此不能丢掉阶乘约定。外积平方中的两个自身平方含重复因子为零，两个交叉项都是$e^1\wedge e^2\wedge e^3\wedge e^4$，因为交换两个二次块的符号为正。故
\[
 \omega\wedge\omega=2e^1\wedge e^2\wedge e^3\wedge e^4,
 \qquad \iota_{e_1}\omega=e^2.
\]
若$\omega=\lambda\wedge\mu$，其平方含两次$\lambda$与$\mu$，由交替性必须为零，与上式矛盾。这个判别已足以否定本例的可分解性，不需要预先建立二形式的全部正规形理论。
\end{proof}

\Needspace{11\baselineskip}
\subsection*{习题6.6：协变张量与形式的拉回}
设$F:\R^2\to\R^2$，$(x,y)=F(u,v)=(u,v+u^2)$。求$F^*(\dd x\otimes\dd y)$和$F^*(\dd x\wedge\dd y)$。令$Z=a\partial_u+b\partial_v$，再计算其对第二个结果的内插。
\begin{proof}[解答]
由$\dd x$和$\dd y$的逐点对偶拉回，$F^*\dd x=\dd u$、$F^*\dd y=\dd v+2u\dd u$。因拉回保持相应的乘法，
\begin{align*}
 F^*(\dd x\otimes\dd y)&=\dd u\otimes\dd v+2u\dd u\otimes\dd u,\\
 F^*(\dd x\wedge\dd y)&=\dd u\wedge\dd v.
\end{align*}
第二式的额外项因$\dd u\wedge\dd u=0$消失，第一式则不能消掉同样的普通张量积。内插的带符号乘法式给出
\[
 \iota_Z(\dd u\wedge\dd v)=a\dd v-b\dd u.
\]
将它作用于$A\partial_u+B\partial_v$得到$aB-bA$，与二形式直接作用于两向量的行列式一致。所有运算逐点且分量光滑，故结果确为光滑张量场或形式。
\end{proof}

\Needspace{12\baselineskip}
\subsection*{习题6.7：Möbius线丛的平方与拉回}
对例\ref{ex:mobius-line}中的$L$，证明$L\otimes L$平凡。再取$f:S^1\to S^1$、$f([t])=[2t]$，证明$f^*L$平凡。说明这些结论为什么不使$L$本身平凡。
\begin{proof}[解答]
在角提升$t$中，令$e_t=[t,1]$为纤维基。提升到$t+2\pi$时$e_{t+2\pi}=-e_t$，所以$e_{t+2\pi}\otimes e_{t+2\pi}=e_t\otimes e_t$。于是局部截面$e_t\otimes e_t$拼成全局光滑无零截面，秩一情形即给出全局标架，所以$L\otimes L$平凡。也可从过渡函数平方为一直接看出。
对拉回丛定义$s([t])=([t],[2t,1])$。若将$t$改为$t+2\pi$，第二分量变为$[2t+4\pi,1]=[2t,1]$，两次反号相消，故定义与代表无关；局部表达光滑且处处非零，所以给出$f^*L$的全局标架。两种构造都消除了反号拼接，而原来的$L$仍受反周期系数的中间值障碍限制。不同丛的平凡性不能相互替代。
\end{proof}

\Needspace{12\baselineskip}
\subsection*{习题6.8：常秩与跳秩的丛映射}
在底空间$\R$的两个平凡秩二丛之间，分别给出矩阵
\[
 B(x)=\begin{pmatrix}1&x\\0&0\end{pmatrix},\qquad
 A(x)=\begin{pmatrix}x&0\\0&1\end{pmatrix}.
\]
求核与像，检验它们何时构成光滑子丛。
\begin{proof}[解答]
$B$恒秩一，核由光滑无零截面$(-x,1)$张成，像由$(1,0)$张成，所以核与像均为平凡秩一子丛。这也直接展示了块消元证明中局部标架的含义。
对$A$，若$x\ne0$，映射可逆，核为零，像为$\R^2$；在$x=0$，核为$\operatorname{span}(1,0)$，像为$\operatorname{span}(0,1)$。纤维维数在原点跳跃，不满足原点邻域上的固定秩局部平凡化要求，故点态核与点态像的整体都不是该邻域上的向量子丛。即使各纤维逐点都是线性子空间，也不足以定义光滑子丛。
\end{proof}

